\documentclass[11pt]{article}
\usepackage[T1]{fontenc}
\usepackage[margin=0.9in]{geometry}
\usepackage{amsmath,amssymb,booktabs,array,longtable}
\usepackage[hidelinks]{hyperref}
\newcommand{\Normal}{\mathcal N}
\newcommand{\BF}{\mathrm{BF}}
\newcommand{\E}{\mathbb E}
\newcommand{\KL}{\operatorname{KL}}
\newcommand{\ERSS}{\operatorname{ERSS}}
\newcommand{\ind}{\mathbf 1}
\newcommand{\T}{^{\mathsf T}}
\newcommand{\yb}{\mathbf y_c}
\newcommand{\etal}{\overline{\boldsymbol\eta}}
\newcommand{\pkg}{\texttt{susieRSlidePrior}}
\setlength{\parskip}{0.35em}
\setlength{\parindent}{0pt}
\setlength{\emergencystretch}{2em}
\allowdisplaybreaks
\title{SuSiE with a fixed finite slider prior\\
\large Model, IBSS implementation, and optional outer empirical Bayes}
\author{Method specification for \pkg{} 0.3.0}
\date{27 September 2026}
\begin{document}
\maketitle

\section{Scope and implemented model}
\label{sec:model}
The current fitter, \texttt{susieRSlidePrior::susie()}, extends the
slider coding in the supplied manuscript~\cite{slider} within the SuSiE
framework~\cite{susie}. It integrates a Gaussian effect size and a finite
slider distribution. \textbf{The slider support and its prior probabilities
are fixed throughout a call to the fitter.} Learning those probabilities is
an optional external procedure, derived separately in Section~\ref{sec:outer}.
Learning the Gaussian variances is a different option and does not learn
the slider probabilities.

\subsection{Genotypes, coding, centering and scaling}
Let $X=(x_{ij})$ contain $n$ samples and $p$ SNPs. The slider interface accepts
complete numeric hard-call genotypes $x_{ij}\in\{0,1,2\}$, including sparse
matrices with those values. Dosages and missing genotypes are rejected.
Define $h_{ij}=\ind\{x_{ij}=1\}$ and
\begin{equation}
 f_\delta(x)=x+\delta\ind\{x=1\}.
 \label{eq:coding}
\end{equation}
The values $\delta=-1,0,1$ give recessive, additive and dominant coding
relative to the counted allele, respectively.

With an intercept, let $\bar x_j$, $\bar h_j$ and $\bar y$ be sample means,
and define
\begin{equation}
 \yb=\mathbf y-\bar y\mathbf1,\qquad
 \mathbf a_j=\frac{\mathbf x_j-\bar x_j\mathbf1}{s_j},\qquad
 \mathbf h_j=\frac{\mathbf h_j^{\rm raw}-\bar h_j\mathbf1}{s_j},\qquad
 \mathbf z_{jk}=\mathbf a_j+d_k\mathbf h_j.
 \label{eq:design}
\end{equation}
Under \texttt{standardize=TRUE}, $s_j$ is the original additive-genotype
sample standard deviation; otherwise $s_j=1$. The engine replaces a zero
scale by one. A constant centered column therefore contains no signal
information. The scale stays the same for every slider value: transformed
columns are not standardized separately. With \texttt{intercept=FALSE},
omit centering, set the mean offsets to zero, and omit the intercept below.
The option \texttt{na.rm=TRUE} can remove observations with missing
phenotypes; genotype completeness is still required.

\subsection{A fixed 17-point prior and the genotype-count rule}
The default grid is
\begin{equation}
 K=17,\qquad d_k=-1+\frac{k}{8},\quad k=0,\ldots,16,
 \qquad w_k=\frac1{17}.
 \label{eq:grid}
\end{equation}
Its spacing is $0.125$, with equal mass at the endpoints and all other
points, including $-0.5,0,0.5$. This is an exact finite distribution, not
quadrature for a continuous prior. User-supplied nonnegative weights are
normalized to sum to one and then held fixed; zero weights stay zero.
The interface also permits other strictly increasing grids within
$[-1,1]$ that include zero. Equations below use the default 17-point indexing.

By default, \texttt{min\_obs=5}. Let $\mathcal J_{\rm free}$ contain the
SNPs with at least \texttt{min\_obs} observations in \emph{each} genotype
class. Counts use the retained observations and uncentered genotypes.
The effective prior is
\begin{equation}
 w_{jk}=\Pr(C_\ell=k\mid J_\ell=j)=
 \begin{cases}
 w_k,&j\in\mathcal J_{\rm free},\\
 \ind\{k=8\},&j\notin\mathcal J_{\rm free}.
 \end{cases}
 \label{eq:forced-prior}
\end{equation}
Thus a count-forced SNP uses additive coding even if the supplied $w_8=0$.
Its additive Bayes factor is not multiplied by $w_8$.
Setting \texttt{min\_obs=0} disables the count restriction.

\subsection{Single-effect components and likelihood convention}
Independently for $\ell=1,\ldots,L$, let
\begin{align}
 J_\ell&\sim\operatorname{Categorical}(\pi_1,\ldots,\pi_p),&
 C_\ell\mid J_\ell=j&\sim\operatorname{Categorical}(w_{j0},\ldots,w_{j16}),
 \nonumber\\
 \delta_\ell&=d_{C_\ell},&\beta_\ell&\sim\Normal(0,V_\ell),\nonumber\\
 \boldsymbol\eta_\ell&=\beta_\ell\mathbf z_{J_\ell C_\ell},&
 \mathbf y&=a_0\mathbf1+\sum_{\ell=1}^L\boldsymbol\eta_\ell+
 \boldsymbol\epsilon,\quad\boldsymbol\epsilon\sim\Normal(0,\sigma^2I_n).
 \label{eq:model}
\end{align}
The Gaussian coefficient is independent of SNP and slider a priori.
The fixed SNP probabilities $\pi_j$ usually equal $1/p$. For free SNPs,
the slider prior is independent of SNP identity. Each component selects
one SNP and one slider, not $p$ independent slider draws. Several
components may select the same SNP with different sliders.

The centered design makes the fitted intercept in these coordinates
$\widehat a_0=\bar y$. The package evaluates the full $n$-observation
Gaussian likelihood at this fitted intercept. It therefore retains $n$
in the likelihood normalizer and residual-variance update. Although
$\yb$ is centered, we do not declare it to have a nonsingular
$n$-dimensional centered-data distribution: equation~\eqref{eq:model}
defines the likelihood before fitting the intercept. The alternative
intercept-integrated convention is described in Section~\ref{sec:intercept}.

The main formulas assume $V_\ell>0$, $\sigma^2>0$ and no explicit null
column. Section~\ref{sec:zero} gives the implementation's null conventions.

\section{Exact single-effect posterior for a fixed residual}
\label{sec:ser}
For the residual vector $\mathbf r_\ell$ used to update component $\ell$,
compute
\begin{align}
 A_j&=\mathbf a_j\T\mathbf a_j,&
 B_j&=\mathbf a_j\T\mathbf h_j,&
 D_j&=\mathbf h_j\T\mathbf h_j,\nonumber\\
 t_{\ell j}&=\mathbf a_j\T\mathbf r_\ell,&
 u_{\ell j}&=\mathbf h_j\T\mathbf r_\ell.
 \label{eq:stats}
\end{align}
Precompute $A_j,B_j,D_j$ once. The two residual scores must be refreshed
when the residual changes. For every grid point,
\begin{equation}
 S_{jk}=A_j+2d_kB_j+d_k^2D_j,\qquad
 T_{\ell jk}=t_{\ell j}+d_ku_{\ell j}.
 \label{eq:ST}
\end{equation}
No $n\times17p$ design is constructed.

\subsection{Gaussian integration and joint posterior weights}
Conditionally on SNP $j$ and grid point $k$, the coefficient posterior is
Gaussian with
\begin{align}
 v_{\ell jk}&=\frac{V_\ell\sigma^2}{\sigma^2+V_\ell S_{jk}},&
 m_{\ell jk}&=\frac{V_\ell T_{\ell jk}}{\sigma^2+V_\ell S_{jk}},
 \label{eq:moments}\\
 b_{\ell jk}:=\log\BF_{\ell jk}
 &=-\frac12\log\left(1+\frac{V_\ell S_{jk}}{\sigma^2}\right)
 +\frac{V_\ell T_{\ell jk}^2}
 {2\sigma^2(\sigma^2+V_\ell S_{jk})}.
 \label{eq:bf}
\end{align}
The Bayes factor compares that candidate with a zero-effect regression
at the same noise variance. If $S_{jk}=0$, the design is zero, so
$T_{\ell jk}=0$, $b_{\ell jk}=0$, $m_{\ell jk}=0$ and $v_{\ell jk}=V_\ell$.

The single-effect evidence relative to its null density $p_0$ and the
posterior are
\begin{align}
 Z_\ell&=\sum_j\sum_k\pi_jw_{jk}\exp(b_{\ell jk}),&
 p(\mathbf r_\ell)&=p_0(\mathbf r_\ell)Z_\ell,\label{eq:evidence}\\
 \rho_{\ell jk}&=\frac{\pi_jw_{jk}\exp(b_{\ell jk})}{Z_\ell},&
 \beta_\ell\mid j,k,\mathbf r_\ell&\sim\Normal(m_{\ell jk},v_{\ell jk}).
 \label{eq:rho}
\end{align}
Normalize $\log\pi_j+\log w_{jk}+b_{\ell jk}$ over all pairs by
log-sum-exp, interpreting zero prior masses as log mass $-\infty$.
The SNP evidence averages Bayes factors over the slider prior; it does
not maximize over grid points.

Define
\begin{equation}
 \alpha_{\ell j}=\sum_k\rho_{\ell jk},\qquad
 \omega_{\ell jk}=\rho_{\ell jk}/\alpha_{\ell j}
 \quad\text{when }\alpha_{\ell j}>0.
 \label{eq:conditional}
\end{equation}
Conditional moments for zero-prior pairs are only storage placeholders;
inference uses pairs with positive posterior mass.
The conditional coefficient distribution given a SNP is a finite
Gaussian mixture. A single Gaussian with matched moments generally
does not reproduce its sign probabilities or posterior samples.

\subsection{Moments used by the implementation}
Write $M^{(2)}_{\ell jk}=v_{\ell jk}+m_{\ell jk}^2$. The conditional
moments stored for each SNP and component are
\begin{align}
 \mu_{\ell j}&=\sum_k\omega_{\ell jk}m_{\ell jk},&
 \mu^{(2)}_{\ell j}&=\sum_k\omega_{\ell jk}M^{(2)}_{\ell jk},\nonumber\\
 \bar\delta_{\ell j}&=\sum_k\omega_{\ell jk}d_k,&
 c_{\ell j}&=\sum_k\omega_{\ell jk}d_km_{\ell jk},\nonumber\\
 e_{\ell j}&=\sum_k\omega_{\ell jk}d_kM^{(2)}_{\ell jk},&
 f_{\ell j}&=\sum_k\omega_{\ell jk}d_k^2M^{(2)}_{\ell jk}.
 \label{eq:aggregates}
\end{align}
Here $c=\E(\beta\delta\mid j)$, $e=\E(\beta^2\delta\mid j)$ and
$f=\E(\beta^2\delta^2\mid j)$. A zero-probability SNP contributes zero
to joint moments regardless of the conditional values used internally.

Let $X_c,H_c$ have columns $\mathbf a_j,\mathbf h_j$. The component mean is
\begin{equation}
 \etal_\ell=X_c\mathbf a_\ell^*+H_c\mathbf h_\ell^*,\qquad
 a^*_{\ell j}=\alpha_{\ell j}\mu_{\ell j},\quad
 h^*_{\ell j}=\alpha_{\ell j}c_{\ell j}.
 \label{eq:fit}
\end{equation}
The joint moment $c_{\ell j}$ cannot in general be replaced by
$\mu_{\ell j}\bar\delta_{\ell j}$.

\section{Fixed-prior IBSS and the variance-update schedule}
\label{sec:ibss}
Use the SuSiE variational factorization
\begin{equation}
 q(J,C,\beta)=\prod_{\ell=1}^Lq_\ell(J_\ell,C_\ell,\beta_\ell).
 \label{eq:q}
\end{equation}
Dependence between SNP, slider and coefficient is retained within a
component; independence is imposed only across components. For $L=1$
the fixed-hyperparameter posterior is exact. For $L>1$ it is variational.

\subsection{One component update}
Maintain the total fitted vector $\sum_b\etal_b$. For component $\ell$,
restore its old contribution and form
\begin{equation}
 \mathbf r_\ell=\yb-\sum_{b\ne\ell}\etal_b.
 \label{eq:residual}
\end{equation}
Compute the two scores in equation~\eqref{eq:stats}; update the prior
variance according to the selected policy below; evaluate the exact grid
posterior; and replace the component fit using equation~\eqref{eq:fit}.
Use this updated fit immediately in the next component's residual.
The noise variance $\sigma^2$ and all $w_k$ stay fixed during the sweep.

\subsection{Effect-prior variance $V_\ell$: three relevant policies}
\begin{enumerate}
 \item \textbf{Fixed:} \texttt{estimate\_prior\_variance=FALSE} keeps
 $V_\ell$ fixed during fitting. No numerical prior-variance optimization
 runs, regardless of the value supplied for \texttt{estimate\_prior\_method}.
 \item \textbf{Component-level EM:}
 \texttt{estimate\_prior\_variance=TRUE} and
 \texttt{estimate\_prior\_method="EM"} first compute the posterior at the
 current $V_\ell$, then set
 \begin{equation}
 V_\ell^{\rm new}=\sum_{j,k}\rho_{\ell jk}
       \bigl(v_{\ell jk}+m_{\ell jk}^2\bigr)
       =\sum_j\alpha_{\ell j}\mu^{(2)}_{\ell j}.
 \label{eq:Vupdate}
 \end{equation}
 The right-hand side uses the pre-update posterior. The slider code
 then evaluates the grid posterior and KL again at the new $V_\ell$,
 using the same residual scores, before updating the fitted contribution.
 This is one moment update and one posterior refresh per component,
 with no inner numerical optimization loop.
 \item \textbf{Numerical empirical Bayes, the current default:}
 \texttt{estimate\_prior\_method="optim"} uses bounded Brent optimization
 of $\log Z_\ell(V)$ before storing the posterior. Each trial $V$ evaluates
 the finite mixture again, but reuses the residual scores. The engine
 retains the previous $V$ if the candidate reduces the SER evidence.
 It also compares with $V=0$ using \texttt{check\_null\_threshold}
 (default zero). A positive threshold deliberately favors a null effect
 and need not preserve unpenalized-ELBO ascent.
\end{enumerate}
The additional \texttt{"simple"} option retains the current positive
variance unless the same null comparison selects zero; it is not the
moment EM update. None of these options updates the $w_k$.

\subsection{One complete sweep and stopping}
The implemented order is:
\begin{enumerate}
 \item Initialize positive noise variance, effect-prior variances and
 zero component means; compute genotype statistics and count restrictions.
 \item Holding $\sigma^2$ and $\mathbf w$ fixed, update components
 $1,\ldots,L$ in order, including any component-level $V_\ell$ update.
 \item Evaluate the ELBO and check convergence. With the normal finite
 ELBO path, after the first sweep convergence requires a nonnegative
 ELBO increase smaller than \texttt{tol}.
 \item If converged, stop. Otherwise, if
 \texttt{estimate\_residual\_variance=TRUE}, update $\sigma^2$ by
 equation~\eqref{eq:sigma}, apply its configured bounds, and start
 the next sweep. If that option is false, keep $\sigma^2$ fixed.
\end{enumerate}
The fitter stops after at most \texttt{max\_iter} sweeps and warns when
it has not converged. A nonconverged last iteration can already have
updated $\sigma^2$; its stored posterior and ELBO should not be interpreted
as a converged fit at that final variance. Final null-effect trimming is
described in Section~\ref{sec:zero}.

Thus only the \emph{residual variance} is necessarily fixed over a whole
sweep. When enabled, estimation of the \emph{effect-prior variance} occurs
inside component updates. A scheme that batches all $V_\ell$ EM updates
after a sweep is another possible algorithm, not the current schedule.

\section{Residual variance, ELBO and correctness of the updates}
\label{sec:objective}
The expected squared residual is
\begin{align}
 R&=\left\|\yb-\sum_\ell\etal_\ell\right\|^2,\nonumber\\
 \ERSS&=R+\sum_\ell\left\{
  \sum_{j,k}\rho_{\ell jk}(v_{\ell jk}+m_{\ell jk}^2)S_{jk}
                        -\|\etal_\ell\|^2\right\}.
 \label{eq:erss}
\end{align}
The code evaluates the inner second moment equivalently as
\begin{equation}
 \sum_j\alpha_{\ell j}
  \left(A_j\mu^{(2)}_{\ell j}+2B_je_{\ell j}+D_jf_{\ell j}\right).
 \label{eq:erss-aggregate}
\end{equation}
This retains uncertainty in SNP, slider and coefficient. Replacing it
by a squared posterior-mean prediction underestimates the expected loss.
The Gaussian residual-variance update used by this interface is
\begin{equation}
 \boxed{(\sigma^2)^{\rm new}=\ERSS/n.}
 \label{eq:sigma}
\end{equation}
It uses stored moments and component fits. It requires neither a new
set of sample-level residual scores nor a 17-point likelihood search.
The accepted residual-method labels \texttt{"MoM"} and \texttt{"MLE"}
both reach this update in the supported ordinary Gaussian slider model;
the NIG residual model is not supported by this interface.

For positive-variance components, the ELBO at the fitted intercept is
\begin{align}
 \mathcal F&=-\frac n2\log(2\pi\sigma^2)-\frac{\ERSS}{2\sigma^2}
             -\sum_\ell K_\ell,\label{eq:elbo}\\
 K_\ell&=\sum_{j,k}\rho_{\ell jk}
                 \log\frac{\rho_{\ell jk}}{\pi_jw_{jk}}\nonumber\\
 &\quad+\frac12\sum_{j,k}\rho_{\ell jk}
 \left\{\frac{v_{\ell jk}+m_{\ell jk}^2}{V_\ell}-1+
                    \log\frac{V_\ell}{v_{\ell jk}}\right\}.
 \label{eq:kl}
\end{align}
Zero-mass terms contribute zero, and posterior mass cannot occur outside
prior support. At the exact SER update the code can compute the same KL
using the evidence identity
\begin{equation}
 K_\ell=-\log Z_\ell+
 \frac{2\mathbf r_\ell\T\etal_\ell-
      \E_{q_\ell}\|\boldsymbol\eta_\ell\|^2}{2\sigma^2}.
 \label{eq:kl-short}
\end{equation}
This identity applies to the residual and parameters used for that SER
update. The stored KL itself remains valid as other components change.

For fixed hyperparameters and other factors,
\begin{equation}
 \E_q\left\|\yb-\sum_b\boldsymbol\eta_b\right\|^2
 =\E_{q_\ell}\|\mathbf r_\ell-\boldsymbol\eta_\ell\|^2+C_{-\ell}.
 \label{eq:residual-identity}
\end{equation}
Consequently, the coordinate objective is a constant plus
$\log p(\mathbf r_\ell)-\KL(q_\ell\|p_{\rm SER,\ell})$.
The exact grid posterior maximizes it. The EM update~\eqref{eq:Vupdate}
and residual update~\eqref{eq:sigma} maximize their parameter blocks with
the appropriate posterior fixed; the additional posterior refresh is
another coordinate ascent step. These ideal updates are nondecreasing
in the ELBO, without a global-optimum guarantee. Numerical tolerances,
positive null thresholds and final threshold-based trimming are separate
implementation choices and should not be included in an unconditional
monotonicity claim.

\subsection{Alternative intercept-integrated convention}
\label{sec:intercept}
The supplied manuscript also uses coordinates $U\T\mathbf y$ with
$U\T U=I_{n-1}$ and $UU\T=I_n-\mathbf1\mathbf1\T/n$.
Integrating out an intercept with a flat prior gives an effective
dimension $N=n-1$. Under that convention, replace $n$ by $N$ in
equations~\eqref{eq:sigma} and~\eqref{eq:elbo}, and use the corresponding
projected vectors. All centered inner products and SER updates at fixed
$\sigma^2$ coincide. \textbf{The current package uses $n$, not $n-1$,}
so the main specification above follows that convention. With no
intercept both conventions use $n$.

\section{Posterior summaries and special cases}
\label{sec:summaries}
\subsection{PIPs, slider summaries and prediction}
For the components retained as active by the package,
\begin{equation}
 \operatorname{PIP}_j=1-\prod_{\ell\in\mathcal L_{\rm active}}
                          (1-\alpha_{\ell j}).
 \label{eq:pip}
\end{equation}
Slider summaries are component-specific:
$\E_q(\delta_\ell\mid J_\ell=j)=\bar\delta_{\ell j}$ and
$\Pr_q(\delta_\ell=0\mid J_\ell=j)=\omega_{\ell j8}$.
Several effects selecting the same SNP do not imply a unique SNP-wide
slider. If all prior mass is placed at zero, the model reduces to additive
SuSiE under matching hyperparameters, scaling and update policies.

On the original genotype scale, the posterior-mean coefficients and
intercept are
\begin{align}
 b_j^x&=s_j^{-1}\sum_\ell\alpha_{\ell j}\mu_{\ell j},&
 b_j^h&=s_j^{-1}\sum_\ell\alpha_{\ell j}c_{\ell j},\nonumber\\
 b_0&=\bar y-\sum_j\bar x_j b_j^x-\sum_j\bar h_j b_j^h,\nonumber\\
 \widehat y_i&=b_0+\sum_jx_{ij}b_j^x+\sum_j\ind\{x_{ij}=1\}b_j^h.
 \label{eq:prediction}
\end{align}
Without an intercept, $b_0=0$. New genotypes must use the same SNP order
and allele orientation. Posterior sign probabilities and coefficient
sampling use the Gaussian mixture over the grid, preserving dependence
between slider and coefficient.

\subsection{Credible sets and their purity diagnostic}
For each eligible component, sort SNPs by $\alpha_{\ell j}$ and take the
smallest prefix attaining the requested coverage (default $0.95$).
Probabilities are marginalized over the slider before constructing a
set. Sets containing the explicit null candidate are not reported;
duplicate SNP sets are removed.

For purity, the implementation forms a component-specific design with
columns $\mathbf x_j+\bar\delta_{\ell j}\mathbf h_j^{\rm raw}$ and computes
absolute correlations among its selected SNPs. The default minimum
absolute correlation threshold is $0.5$. Large sets may be subsampled
for this calculation according to \texttt{n\_purity}; optional median
correlation filtering is also supported. This is a diagnostic based on
posterior-mean coding, not a correlation integrated over slider uncertainty.
The credible-set probability calculation still uses the marginalized
$\alpha$. Setting \texttt{coverage=NULL} skips set construction and purity.

\subsection{Zero effects, null candidates and numerical trimming}
\label{sec:zero}
At $V_\ell=0$, the coefficient is deterministically zero; use the limiting
formulas $m=v=0$ and $\BF=1$ rather than divide by $V_\ell$ in
equation~\eqref{eq:kl}. Such a component has no fitted contribution and
no information about inheritance mode. Its stored categorical posterior
can still equal the prior; that is bookkeeping, not an observed slider.

At finalization the engine sets $V_\ell<\texttt{prior\_tol}$ to zero
(default tolerance $10^{-9}$), resets its coefficient moments and fits,
and excludes zero-variance components from slider-prior learning counts.
Thresholding a small positive variance is a numerical convention, not
the exact positive-variance model used in the main derivation.

With \texttt{null\_weight}>0, an additional zero-design candidate has prior
probability $\pi_0$. Its Bayes factor is one, so
\begin{equation}
 Z_\ell=\pi_0+\sum_{j=1}^p\sum_k\pi_jw_{jk}\BF_{\ell jk},\qquad
 \rho_{\ell0}=\pi_0/Z_\ell,
 \label{eq:null}
\end{equation}
where the SNP weights now sum to $1-\pi_0$. This candidate contributes zero
to predictions and ERSS and is excluded from slider-learning counts.
Internally it can use additive coding as a placeholder; it is not evidence
for an additive biological effect. If the latent Gaussian coefficient is
retained for this zero-design candidate, its posterior is its prior and
its second moment is $V_\ell$, as used by the coefficient-variance update.

\section{Using the package and interpreting its output}
\label{sec:api}
This call keeps the slider prior and coefficient-prior variance fixed
while estimating the residual variance after each non-final sweep:
\begin{verbatim}
fit <- susieRSlidePrior::susie(
  X, y, L = 10,
  delta_grid = seq(-1, 1, length.out = 17),
  delta_prior = rep(1/17, 17),
  estimate_prior_variance = FALSE,
  estimate_residual_variance = TRUE
)
\end{verbatim}
Unless supplied otherwise, the initial $V_\ell$ is
$0.2\,\operatorname{var}(y)$, determined by
\texttt{scaled\_prior\_variance=0.2}, and the initial noise variance is
$\operatorname{var}(y)$. To request an absolute $V_\ell=V_0$, supply
\texttt{scaled\_prior\_variance=V0/var(y)} on the retained data.
The interface defaults to at most ten effects, \texttt{max\_iter=100},
\texttt{tol=1e-3}, and estimation of both Gaussian variances with
\texttt{estimate\_prior\_method="optim"}. The explicit call above changes
the coefficient-variance policy while retaining a fixed slider prior.

For a moment-based coefficient-variance update, instead use
\begin{verbatim}
fit_em <- susieRSlidePrior::susie(
  X, y, L = 10,
  delta_prior = rep(1/17, 17),
  estimate_prior_variance = TRUE,
  estimate_prior_method = "EM",
  estimate_residual_variance = TRUE
)
\end{verbatim}
Here \texttt{EM} refers to the Gaussian coefficient variance, not the slider
prior probabilities. The exact support and supplied probabilities are
returned in \texttt{delta\_grid} and \texttt{delta\_prior}.

In the following table the R array index for mathematical grid point
$k=0,\ldots,16$ is $k+1$. The dimensions are $L\times p$ or $L\times p\times K$;
an explicit null column, if present, adds an internal candidate.
{\small
\begin{longtable}{@{}p{0.36\textwidth}p{0.59\textwidth}@{}}
\toprule
Returned field & Meaning\\
\midrule
\endhead
\path{alpha} & $\alpha_{\ell j}$, the component SNP probabilities.\\
\path{alpha_delta} & $\rho_{\ell jk}$, joint SNP--slider probabilities;
 summing over the third dimension gives \texttt{alpha}.\\
\path{mu_grid}, \path{mu2_grid} & $m_{\ell jk}$ and
 $v_{\ell jk}+m_{\ell jk}^2$, conditionally on a SNP--slider pair.\\
\path{mu}, \path{mu2} & $\mu_{\ell j}$ and $\mu^{(2)}_{\ell j}$.\\
\path{delta} & $\bar\delta_{\ell j}$, the posterior mean slider
 conditional on the SNP and component.\\
\path{mu_delta} & $c_{\ell j}=\E_q(\beta_\ell\delta_\ell\mid J_\ell=j)$.\\
\path{mu2_delta} & $e_{\ell j}=\E_q(\beta_\ell^2\delta_\ell\mid J_\ell=j)$.\\
\path{mu2_delta2} & $f_{\ell j}=\E_q(\beta_\ell^2\delta_\ell^2\mid J_\ell=j)$.\\
\path{delta_weights} & $\sum_j\rho_{\ell jk}$ for each component and grid point.\\
\path{delta_prior_counts} & Component-by-grid counts excluding count-forced
 SNPs, the explicit null candidate and zero-variance components. No
 credible-set quality filter is applied.\\
\path{expected_squared_residuals} & The mixture-correct $\ERSS$.\\
\path{V}, \path{sigma2}, \path{elbo} & Gaussian prior variances, residual
 variance, and the objective trace used by IBSS.\\
\path{delta_forced}, \path{genotype_counts} & The count restriction and
 the observed numbers of genotype classes 0, 1 and 2.\\
\path{pip}, \path{sets} & SNP PIPs and credible-set output.\\
\path{delta_cs} & A per-member summary table and a CS-by-SNP matrix of
 posterior mean sliders, indexed by the reported components.\\
\bottomrule
\end{longtable}
}

\texttt{coef(fit)} returns the additive and heterozygote coefficients
in equation~\eqref{eq:prediction}, including the intercept row.
\texttt{predict(fit,newx=Xnew)} uses the original hard-call coding.
\texttt{susieRSlidePrior::slider\_cs\_table(fit)} gives the credible-set
member table. The posterior sampling helper returns additive and
heterozygote coefficient draws jointly.

A warm start is supplied through \texttt{model\_init=fit}. It requires
compatible grid, SNP ordering, dimensions, scaling and count restrictions.
It copies posterior state and effect variances while retaining newly
supplied SNP and slider priors. This permits an external procedure to
refit with changed \texttt{delta\_prior}.

For comparison, explicitly setting \texttt{delta} uses fixed coding and
overrides the finite prior. Setting \texttt{delta\_prior=NULL} with
\texttt{delta=NULL} selects the legacy continuous plug-in slider.
These are separate fitting modes. The inherited function
\texttt{susieRSlidePrior::susie\_additive()} is explicitly additive;
inherited summary-statistics interfaces do not implement this
individual-level slider model.

\section{Computational cost}
\label{sec:cost}
The comparison below concerns individual-level additive SuSiE with a
Gaussian coefficient prior, the same $n,p,L$, and matching variance
policies. Both methods integrate the Gaussian coefficient analytically.

\subsection{Sample-level work and finite-grid work}
Setup is $O(np)$. The slider additionally computes heterozygote counts
and the sufficient statistics $B_j,D_j$. At each component residual,
additive SuSiE needs $\mathbf a_j\T\mathbf r_\ell$; the slider also needs
$\mathbf h_j\T\mathbf r_\ell$. This is one extra length-$n$ inner product
per SNP, approximately $n$ multiply-accumulate operations, or $2n-1$
separate multiplications and additions in dense arithmetic.

After these scores, the 17 candidates need only scalar calculations
through equation~\eqref{eq:ST}. They do not require 17 passes over samples.
A full component update also adds one multiplication by $H_c$ when
refreshing the fitted vector.
\begin{center}
\small
\begin{tabular}{@{}p{0.25\textwidth}p{0.29\textwidth}p{0.38\textwidth}@{}}
\toprule
Work per component & Additive Gaussian & Finite slider\\
\midrule
Residual scores & $X_c\T\mathbf r_\ell$ &
 $X_c\T\mathbf r_\ell,\ H_c\T\mathbf r_\ell$\\[0.3em]
Candidate evaluation & $p$ candidates & $Kp$ pairs, $K=17$\\[0.3em]
Fitted contribution & $X_c\mathbf a_\ell^*$ &
 $X_c\mathbf a_\ell^*+H_c\mathbf h_\ell^*$\\
\bottomrule
\end{tabular}
\end{center}
The leading dense matrix-vector work is approximately $4np$ operations
per additive component update and $8np$ per slider component update,
counting multiplication and addition separately. Thus the additional
leading matrix work is about $4np$ per component or $4nLp$ per sweep,
plus $O(KLp)$ scalar work. Binary group sums can lower arithmetic costs.

For $T$ sweeps, with a bounded number of candidate evaluations per update,
the mathematical algorithm can be organized as
\begin{equation}
 O\{np+TL(np+Kp)\},\qquad K=17.
 \label{eq:cost}
\end{equation}
The corresponding additive order is $O\{np+TL(np+p)\}$. This is an
arithmetic-complexity comparison, not a measured runtime guarantee or
an exact accounting of R allocations and copying.

\subsection{Why the variance policy matters}
The residual-variance update uses equation~\eqref{eq:erss-aggregate}
and stored fits, taking $O(Lp+nL)$ additional aggregation work in the
current representation. It adds no residual cross-product or grid
likelihood search. The same kind of expected-residual calculation is
also needed for the ELBO.

For the coefficient variance, one EM moment update is $O(p)$ after
aggregation; the implementation's posterior refresh adds one $O(Kp)$
candidate evaluation without repeating sample-level scores. With
\texttt{"optim"}, if $F_{\ell t}$ objective/posterior evaluations are
made for component $\ell$ in sweep $t$, candidate work is instead
\begin{equation}
 O\left\{np+\sum_{t=1}^T\sum_{\ell=1}^L
              \bigl(np+F_{\ell t}Kp\bigr)\right\}.
 \label{eq:optim-cost}
\end{equation}
The present native routine calculates mixture moments as well as evidence
on these calls. Consequently, repeated numerical optimization can be
substantially more expensive than estimating only $\sigma^2$ after a sweep.
Cheap individual EM updates do not guarantee fewer convergence sweeps.

Against the original plug-in slider, both residual scores and both fitted
products already exist. Finite integration replaces its cubic slider
maximization with $K$ Bayes-factor/moment evaluations and mixture sums.
It does not inherently add a further sample-level pass relative to that
slider implementation.

\subsection{Storage and external prior reweighting}
The current fitter retains the joint posterior and conditional grid
moments, requiring $O(KLp)$ storage, in addition to the genotype matrix,
aggregate moments and component fits. A memory-reduced implementation
could retain aggregates and regenerate pair posteriors, but that is not
the current storage policy. Heterozygote indicators are formed in chunks
by default or cached when requested; an $n\times Kp$ matrix is unnecessary.

For a fixed residual and fixed Gaussian variances, cached grid Bayes
factors can be reweighted in $O(Kp)$. An actual outer prior update normally
requires a new IBSS fit because the component residuals change.
The current external-refit approach does not claim that an entire
outer iteration is only an $O(Kp)$ reweighting.

\section{Measured benchmarks}
\label{sec:bench}
These are measured full-fit times on 27 September 2026, not operation
counts. Each reported median uses five \texttt{microbenchmark} repetitions
after one complete warm-up fit. Implementations run in separate sequential
R processes. Timings include input preparation, IBSS and returned-object
construction, and exclude simulation, package loading, file I/O and
explicit pre-fit garbage collection. Garbage collection during a fit
remains included.

All scenarios use $L=10$, \texttt{tol=1e-6}, \texttt{max\_iter=300},
initial $V_\ell=0.5$ and $\sigma^2=0.64$, and \texttt{coverage=NULL}.
The genotypes are independent Binomial$(2,0.3)$ hard calls. Three causal
SNPs have coefficients $(0.9,-0.7,0.8)$ and sliders $(-0.5,0,0.5)$,
with noise standard deviation $0.8$. Each size uses the same fixed
dataset across methods and settings. Finite-prior fits keep every slider
probability at $1/17$; the original slider uses continuous plug-in estimates.
The slider settings are \texttt{min\_obs=5},
\texttt{chunk\_size=1000}, and no heterozygote cache.

\subsection{First run: fixed variances versus coefficient-variance optimization}
\begin{center}
\small
\begin{tabular}{@{}rrlrrr@{}}
\toprule
$n$ & $p$ & Gaussian variances & SuSiE & Original slider & Finite prior\\
\midrule
500 & 1000 & Both fixed & 0.177 & 1.185 & 1.539\\
500 & 1000 & Both estimated; $V$: \texttt{optim} & 0.276 & 5.257 & 10.433\\
1000 & 3000 & Both fixed & 1.288 & 7.125 & 6.211\\
\bottomrule
\end{tabular}
\end{center}
All times are seconds. The corresponding sweep counts, ordered as
SuSiE/original slider/finite prior, are $8/8/7$, $4/8/4$ and $11/8/8$.
The estimated-variance finite-prior fit is about $1.98$ times the original
slider and $37.85$ times standard SuSiE in this small dataset; this
comparison includes repeated numerical optimization of $V_\ell$.
It does not measure only the cost of updating $\sigma^2$.
For the larger fixed-variance case, the interquartile ranges of the
slider times overlap: $7.018$--$7.447$ versus $6.096$--$7.781$ seconds.
The smaller finite-prior median is therefore not an established speed
advantage.

\subsection{Follow-up run: isolate the residual-variance update}
This is a separate paired run at $n=500,p=1000$, with every $V_\ell$
fixed in both settings.
\begin{center}
\small
\begin{tabular}{@{}lrrcc@{}}
\toprule
Implementation & Fixed $\sigma^2$ & Updated $\sigma^2$ &
 \multicolumn{2}{c}{Sweeps}\\
 & (seconds) & (seconds) & Fixed & Updated\\
\midrule
SuSiE & 0.179 & 0.171 & 8 & 7\\
Original slider & 0.813 & 0.970 & 8 & 10\\
Finite prior & 1.058 & 1.200 & 7 & 8\\
\bottomrule
\end{tabular}
\end{center}
For the finite-prior fitter, dividing total time by sweep count gives
approximately $0.151$ and $0.150$ seconds, including amortized setup and
output construction. The total-time difference is consistent with the
additional sweep. This supports the small cost of the residual update.
No component-variance EM timing was measured in these runs.

\subsection{Environment, storage and limits of the comparison}
The implementations are \texttt{susieR} 0.16.6, original
\texttt{susieSlide} 0.2.0 and the finite-prior code 0.3.0.
The first run preceded the package rename, so its recorded finite-prior
package name is \texttt{susieSlide} 0.3.0; the follow-up uses
\pkg{} 0.3.0. The finite-prior build was cleanly compiled with
\texttt{-O2}. The environment was Windows 11, R 4.6.1 and
\texttt{microbenchmark} 1.5.0, with an AMD64 Family 26 Model 112 processor.
OpenMP/MKL/OpenBLAS thread environment variables were set to one.

All warm-up fits converged. These models need not select the same number
of effects or converge in the same number of sweeps. At $500\times1000$,
returned-object sizes were approximately $0.71$, $5.84$ and $10.08$ MiB;
at $1000\times3000$, $2.10$, $27.40$ and $40.34$ MiB, in the same method
order. These are \texttt{object.size} measurements, not peak memory.
The data have no simulated LD, and the repetitions quantify execution
variability on one machine, not variation across loci or signal patterns.
Comparisons of the two variance policies should use the paired
follow-up, rather than mix times from different runs.

The reusable comparison script is
\path{inst/examples/benchmark_slider_prior.R}. Full measurements and
session records are in the analysis workspace under
\path{output/benchmarks/slider_prior_2026-09-27/release/} and
\path{output/benchmarks/slider_prior_2026-09-27/residual_variance_only/}.
The latter run's script is \path{tmp/slider_prior_impl/benchmark_residual_variance.R}.

\section{Optional outer empirical Bayes for the slider probabilities}
\label{sec:outer}
\textbf{This section derives an extension; the current fitter does not
perform these updates internally.} It fits with supplied $\mathbf w$ and
returns posterior counts so an external fit--update--refit workflow can
learn a common prior, in the same organizational style as coding-mixture
prior learning.

\subsection{Unrestricted common-prior update}
First suppose there are no count restrictions, null candidates or
discarded components, and the set of $L$ components is fixed.
At a fixed variational posterior, the part of the objective depending on
$\mathbf w$ is
\begin{equation}
 Q(\mathbf w)=\sum_kN_k\log w_k+\text{constant},\qquad
 N_k=\sum_{\ell=1}^L\sum_j\rho_{\ell jk}.
 \label{eq:Q}
\end{equation}
Since $\sum_kN_k=L$, maximization over the simplex gives
\begin{equation}
 \boxed{w_k^{\rm new}=\frac{N_k}{L}
       =\frac1L\sum_{\ell=1}^L\sum_j\rho_{\ell jk}.}
 \label{eq:wupdate}
\end{equation}
There is one latent slider draw per component, so the denominator is $L$,
not $Lp$ or $KL$. The support $d_k$ never moves. Unlike the present
fixed-prior fit, uniform probabilities here are only the initialization.

\subsection{Count restrictions, nulls and pooling across loci}
With equation~\eqref{eq:forced-prior}, only free SNPs inform $\mathbf w$.
For independent loci $g$ and a specified set of included components
$\mathcal A_g$, define
\begin{align}
 N_k&=\sum_g\sum_{\ell\in\mathcal A_g}
          \sum_{j\in\mathcal J_{{\rm free},g}}\rho_{g\ell jk},&
 D&=\sum_kN_k,\nonumber\\
 w_k^{\rm new}&=N_k/D\quad(D>0).
 \label{eq:restricted-update}
\end{align}
The denominator is the total posterior mass on eligible, non-null SNPs,
not automatically the number of included components. Keep the previous
weights if $D=0$. Each locus retains its own Gaussian variances and
residuals. The package's \texttt{delta\_prior\_counts} supplies the needed
per-component counts with forced SNPs and nulls removed and zero-variance
rows set to zero.

There is no built-in credible-set purity or coverage eligibility filter
on those counts. An external workflow can choose included component rows,
but that choice must be specified. For a fixed model and included set,
equation~\eqref{eq:restricted-update} is the corresponding M-step.
Changing membership according to fitted credible-set quality or a moving
variance threshold is a learning policy; it does not automatically
preserve the full-data EM ascent guarantee. Excluding zero-effect rows
is natural because they contain no slider information, but should be
distinguished from the unrestricted latent-draw counting convention.

\subsection{Outer algorithm and a count-update example}
\begin{enumerate}
 \item Choose the fixed grid, initial probabilities (usually $1/17$),
 participating loci/components, and Gaussian-variance policies.
 \item At fixed $\mathbf w^{(r)}$, fit or warm-start each locus with
 \pkg{}, allowing IBSS to converge to the requested inner tolerance.
 \item Hold those posterior factors fixed. Sum the eligible counts and
 apply equation~\eqref{eq:restricted-update}.
 \item Refit at $\mathbf w^{(r+1)}$. Continue until the chosen weight and
 objective tolerances are met. The final reported posterior must come
 from a fit using the final reported weights.
\end{enumerate}
For one fitted locus, the unfiltered active-row count update is:
\begin{verbatim}
counts <- colSums(fit$delta_prior_counts)
w_new <- if (sum(counts) > 0) counts / sum(counts) else fit$delta_prior
# Refit X and y with the same L, scaling and variance policies:
# delta_prior = w_new, model_init = fit
\end{verbatim}
This snippet is one M-step, not a complete convergence-controlled outer
EM driver. Zero weights cannot revive under unregularized EM; start with
positive mass on all grid values that should remain eligible for learning.

For a fixed model, posterior coordinate updates increase the ELBO and
the count update maximizes it over $\mathbf w$. Alternating these blocks
therefore gives variational EM for $L>1$. With $L=1$ and an exact posterior
E-step, it is ordinary EM. These statements do not assert global
optimality or incorporate additional data-dependent selection rules.

\subsection{Information, regularization and allele orientation}
Seventeen weights have sixteen free parameters. A small number of
signals may provide little information about a population distribution.
For $L=1$, fixed Gaussian variances and an unrestricted common prior,
\begin{equation}
 Z(\mathbf w)=\sum_kw_kB_k,\qquad B_k=\sum_j\pi_j\BF_{jk}.
\end{equation}
The likelihood maximum on the simplex concentrates on a maximizing grid
point (or tied maxima). Pooling loci supplies replicated effects.
An optional Dirichlet prior with parameters $1+\lambda/17$ gives the
penalized update
\begin{equation}
 w_k^{\rm new}=\frac{N_k+\lambda/17}{D+\lambda}.
\end{equation}
This changes the objective to MAP/penalized empirical Bayes and is not
part of the current fixed-prior fitter.

With centered coding or a fitted intercept, reversing the counted allele
maps $(\beta,\delta)$ to $(-\beta,-\delta)$, with an intercept adjustment
on the original scale. Without an intercept, raw-coding allele flips need
not preserve the model. An asymmetric prior requires a consistent allele
convention. In the intercept-adjusted setting, to enforce invariance to allele flips,
impose $w_k=w_{16-k}$. For $k<8$, the unregularized constrained update is
$w_k=w_{16-k}=(N_k+N_{16-k})/(2D)$, while $w_8=N_8/D$.
The fixed uniform prior already satisfies this symmetry. Empirical-Bayes
inference averages slider uncertainty at fitted weights, but does not
integrate uncertainty in the estimated population weights themselves.

\section{Implementation correspondence and verification}
The public interface is in \path{R/slider.R}. The native finite-mixture
kernel is in \path{src/slider_prior.cpp}; its R wrapper and moment helpers
are in \path{R/slider_prior.R}. IBSS dispatch, joint-moment fits, ERSS,
variance hooks and output counts are in \path{R/slider_engine.R}.
Prediction and credible-set summary helpers are in
\path{R/slider_methods.R}. The shared workhorse controls the sweep and
end-of-sweep residual-variance update.

The finite kernel uses stable log-mass normalization, respects zero
prior weights, and handles zero design norms and zero prior variances
by their limiting formulas. It uses higher-precision intermediate
arithmetic before returning ordinary R numeric arrays. Count-forced
SNPs and null candidates use their effective point prior at zero.

Validation includes agreement with direct Gaussian-mixture calculations,
an explicitly expanded 17-coding Gaussian design, point-mass reduction
to fixed coding, mixture-correct moments and predictions, fixed weights
during Gaussian-variance learning, warm starts, and null/count restrictions.
After the package rename, all 567 slider test expectations passed in the
installed release build. This verifies those tested cases; it does not
replace convergence diagnostics for new datasets.

\begin{thebibliography}{9}
\bibitem{slider} W. Denault (September 2026).
\emph{The slider model}. User-supplied manuscript
\texttt{slider\_model-6.pdf}.
\bibitem{susie} G. Wang, A. Sarkar, P. Carbonetto and M. Stephens (2020).
A simple new approach to variable selection in regression, with
application to genetic fine mapping. \emph{JRSS B}, 82, 1273--1300.
\href{https://doi.org/10.1111/rssb.12388}{doi:10.1111/rssb.12388}.
\end{thebibliography}
\end{document}
